\exercisesection{§ 5 的习题}

\begin{enumerate}
\def\labelenumi{\arabic{enumi})}
\item
  确定一个有限二面体群的对称不变元代数（就其 2 维典范表示而言，cf.~§4, no. 2）。
\item
  设 A 是一个主理想环，E 是一个有限秩 l 的自由 A-模，G 是 \(\mathbf{GL}(\mathbf{E})\) 的一个有限子群。假设：
\end{enumerate}

\begin{enumerate}
\def\labelenumi{(\roman{enumi})}
\item
  若 \(q = \operatorname{Card}(G)\)，则元素 \(q.1\)（\(A\) 中的元素）可逆。
\item
  G 由 E 中的伪反射生成（即由使 \((s - 1)(E)\) 是 E 的一个单生成子模的元素 s 生成，cf.~§2, Exerc. 1。）
\end{enumerate}

设 \(\mathbf{S}(\mathrm{E})\) 是 \(\mathrm{E}\) 的对称代数，\(\mathbf{S}(\mathrm{E})^{\mathrm{G}}\) 是 \(\mathbf{S}(\mathrm{E})\) 中由在 \(\mathrm{G}\) 下不变的元素组成的子代数。证明：对从 \(\mathrm{A}\) 到域 \(k\) 中的任何同态，\(\mathbf{S}(\mathrm{E})^{\mathrm{G}} \otimes k\) 可与 \(\mathbf{S}(\mathrm{E} \otimes k)^{\mathrm{G}}\) 等同。由此，应用 Th. 4 推出，\(\mathbf{S}(\mathrm{E})^{\mathrm{G}}\) 是 \(\mathrm{A}\) 上的一个分次多项式代数。

¶ 3) 设 K 是一个特征为零的域，V 是 K 上一个有限维 l 的向量空间，G 是 GL(V) 的一个由伪反射生成的子群。置 q = Card(G)；用 S（相应地，L）表示 V 的对称代数（相应地，外代数）。若 \(x \in V\)，用 x（相应地，\(x'\)）表示它在 S（相应地，L）中的典范像。

\begin{enumerate}
\def\labelenumi{\alph{enumi})}
\item
  设 \(E = S \otimes L\) 是 S 与 L 的张量积。证明 E 上存在唯一的导子 d，使得 \(dx = x'\) 且 \(dx' = 0\) 对所有 \(x \in V\) 成立。
\item
  设 \(S^{G}\) 是 G 的对称不变元代数，\(P_{1},\ldots,P_{l}\) 是 \(S^{G}\) 中使 \(S^{G}=K[P_{1},\ldots,P_{l}]\) 的齐次元素。对每个子集 \(I=\{i_{1},\ldots,i_{r}\}\)（\([1,l]\) 的、满足 \(i_{1}<\cdots<i_{r}\) 的子集），置
\end{enumerate}

\[
\omega_ {\mathrm{I}} = d \mathrm{P} _ {i _ {1}} \dots d \mathrm{P} _ {i _ {r}}.
\]

证明诸 \(\omega_{I}\) 在 S 上线性无关，且都属于 E 中由在 G 下不变的元素组成的子代数 \(E^{G}\)。由此推出，对所有 \(\omega \in E^{G}\)，存在 \(a, c_{I} \in S^{G}\)，其中 \(a \neq 0\)，使得

\[
a \omega = \sum_ {\mathrm{I}} c _ {\mathrm{I}} \omega_ {\mathrm{I}}.
\]

\(c)\) 证明 \(\mathrm{E}^{\mathrm{G}}\) 中包含在 \(S\otimes \bigwedge^{l}\mathrm{V}\) 中的每个元素都形如 \(c.dP_{1}\ldots dP_{l}\)，其中 \(c\in S^{\mathrm{G}}\)。（应用 Prop. 5。）

\begin{enumerate}
\def\labelenumi{\alph{enumi})}
\setcounter{enumi}{3}
\item
  证明诸 \(\omega_{I}\) 构成 \(S^{G}\)-模 \(E^{G}\) 的基。（用 b) 的记号，把关系式 \(a\omega = \sum_{I} c_{I}\omega_{I}\) 的两边乘以一个元素 \(\omega_{J}\)，并应用 c)；由此推出，若 I 是 J 的余集，则 a 整除 \(c_{I}\)；从而 \(^{9}\) 诸 \(\omega_{I}\) 生成 \(E^{G}\)。）
\item
  设 \(S_{n}\)（相应地，\(L_{m}\)）是 \(S\)（相应地，\(L\)）的 \(n\) 次（相应地，\(m\) 次）齐次分量。置
\end{enumerate}

\[
\begin{array}{r l r} & & {\mathrm{E} _ {n, m} = \mathrm{S} _ {n} \otimes \mathrm{L} _ {m}, \quad \mathrm{E} _ {n, m} ^ {\mathrm{G}} = \mathrm{E} ^ {\mathrm{G}} \cap \mathrm{E} _ {n, m}, \quad a _ {n, m} = \dim \mathrm{E} _ {n, m} ^ {\mathrm{G}},} \\ & & {a (\mathrm{X}, \mathrm{Y}) = \sum_ {n, m \geqslant 0} a _ {n, m} \mathrm{X} ^ {n} \mathrm{Y} ^ {m}.} \end{array}
\]

利用 \(d\) )，证明公式

\[
a (\mathrm{X}, \mathrm{Y}) = \prod_ {i = 1} ^ {l} \frac {1 + \mathrm{Y} . \mathrm{X} ^ {p _ {i} - 1}}{1 - \mathrm{X} ^ {p _ {i}}},
\]

其中 \(p_{i} = \deg(P_{i})\)。

\(f)\) 若 \(g \in \mathbf{G}\)，用 \(\operatorname{Tr}_{n,m}(g)\) 表示 \(\mathbf{E}_{n,m}\) 的、由 \(g\) 所定义的自同构的迹；置

\[
\operatorname{Tr} (\mathrm{X}, \mathrm{Y}) (g) = \sum_ {n, m \geqslant 0} \operatorname{Tr} _ {n, m} (g) \mathrm{X} ^ {n} \mathrm{Y} ^ {m}.
\]

证明

\[
\operatorname{Tr} (X, Y) (g) = \frac {\det (1 + Y g)}{\det (1 - X g)}.
\]

\begin{enumerate}
\def\labelenumi{\alph{enumi})}
\setcounter{enumi}{6}
\tightlist
\item
  设 \(q = \operatorname{Card}(\mathbf{G})\)。证明
\end{enumerate}

\[
\frac {1}{q} \sum_ {g \in G} \operatorname{Tr} (X, Y) (g) = a (X, Y),
\]

即

\[
\frac {1}{q} \sum_ {g \in G} \frac {\det (1 + Y g)}{\det (1 - X g)} = \prod_ {i = 1} ^ {l} \frac {1 + Y . X ^ {p _ {i} - 1}}{1 - X ^ {p _ {i}}}.
\]

（利用下列结果：若 \(\mathbf{G}\) 作用在一个有限维向量空间 \(\mathbf{E}\) 上，则 \(\mathbf{E}\) 中在 \(\mathbf{G}\) 下不变的元素所成空间的维数等于 \(\frac{1}{q}\sum_{g\in \mathbf{G}}\mathrm{Tr}_{\mathbf{E}}(g).\)）

当 Y = 0 时这个公式给出什么？

\begin{enumerate}
\def\labelenumi{\alph{enumi})}
\setcounter{enumi}{7}
\tightlist
\item
  对任何整数 \(p \geqslant 0\)，设 \(H_{p}\) 是 G 中以 1 为重数 p 的特征值的元素 \(g \in G\) 组成的集合。~设 \(h_{p} = \text{Card}(H_{p})\)。证明公式
\end{enumerate}

\[
\sum_ {p = 0} ^ {l} h _ {p} \mathbf {T} ^ {p} = \prod_ {i = 1} ^ {l} (p _ {i} - 1 + \mathbf {T}).
\]

（在上面的 \(g\) ) 中的公式里，把 \(Y\) 换成 \(-1 + T(1 - X)\)，然后在所得结果中置 \(X = 1\)。若 \(g \in H_p\)，则项 \(\operatorname{Tr}(X, Y)(g)\) 变成 \(T^p\)。）

¶ 4) *设 \(G _ {1} = \mathbf {S L} _ {3} (\mathbf {F} _ {2})\)。

\begin{enumerate}
\def\labelenumi{\alph{enumi})}
\item
  证明 \(\mathbf{G}_1\) 是一个 168 阶的非交换单群，含有 21 个 2 阶元素。
\item
  证明 \(G_{1}\) 的不可约复表示的次数是 1, 3, 3, 6, 7, 8。
\item
  设 \(\rho: \mathbf{G}_1 \to \mathbf{GL}_3(\mathbf{C})\) 是一个不可约表示 \(^{10}\)，它是 \(\mathbf{G}_1\) 的 3 次表示。若 \(y \in \mathbf{G}_1\) 是 2 阶的，证明 \(\operatorname{Tr}(\rho(y)) = -1\)。由此推出 \(-\rho(y)\) 是一个反射。
\item
  设 G 是 \(\mathbf{GL}_{3}(\mathbf{C})\) 中由诸元素 \(-\rho(y)\)（y 取遍 \(G_{1}\) 中的 2 阶元）生成的子群。证明 G 同构于 \(G_{1} \times \{1, -1\}\)，从而其阶为 336。
\item
  证明特征次数 \(k_{1}, k_{2}, k_{3}\)（\(\mathbf{G}\) 的对称不变元代数的特征次数）等于 4、6 与 14。（利用关系式 \(\prod_{i} k_{i} = 336\) 与 \(\sum_{i} (k_{i} - 1) = 21\)。）
\item
  证明 G 不是一个 Coxeter 群。*
\end{enumerate}

\begin{enumerate}
\def\labelenumi{\arabic{enumi})}
\setcounter{enumi}{4}
\tightlist
\item
  *设 \(K\) 是一个域，\(S = K[X_1, \ldots, X_n]\) 是 \(K\) 上的一个分次多项式代数，由代数无关的齐次元素 \(X_i\)（次数 \(> 0\)）生成。
\end{enumerate}

\begin{enumerate}
\def\labelenumi{\alph{enumi})}
\tightlist
\item
  设 \(Y_{1},\ldots,Y_{n}\) 是 S 中次数 \textgreater0 的齐次元素，\(R = K[Y_{1},\ldots,Y_{n}]\) 是 S 中由这些元素生成的子代数。证明下列性质等价：
\end{enumerate}

\begin{enumerate}
\def\labelenumi{(\roman{enumi})}
\item
  \((\mathrm{Y}_1,\dots ,\mathrm{Y}_n)\) 是一个 S-正则序列。
\item
  S 在 R 上是整的。
\item
  \(S\) 中由 \(Y_{1},\ldots ,Y_{n}\) 生成的理想在 \(S\) 中有有限余维数。
\item
  对任何扩张 \(\overline{\mathbf{K}}\)（\(\mathbf{K}\) 的扩张），方程组 \(Y_{i}(x_{1},\ldots ,x_{n}) = 0\)（\(1\leqslant i\leqslant n\)，\(x_{i}\in \overline{\mathbf{K}}\)）只有平凡解 \((0,\dots ,0)\)。
\end{enumerate}

\(((\mathrm{i}) \Longleftrightarrow (\mathrm{iii})\) 由 Macaulay 的一个定理得出（Commutative Algebra）；\((\mathrm{iii}) \Longleftrightarrow (\mathrm{iv})\) 由零点定理得出（Commutative Algebra, Chap. V, §3, no. 3, Prop. 2）；(ii) \(\Longleftrightarrow (\mathrm{iii})\) 是容易的。）

若这些性质满足，证明诸 \(Y_{i}\) 在 K 上代数无关，且 S 是一个自由 R-模，其秩等于

\[
\frac {\prod_ {i} \deg (\mathrm{Y} _ {i})}{\prod_ {i} \deg (\mathrm{X} _ {i})}.
\]

\begin{enumerate}
\def\labelenumi{\alph{enumi})}
\setcounter{enumi}{1}
\tightlist
\item
  设 G 是分次代数 S 的一个有限自同构群，\(S^{G}\) 是它的不变元子代数，\(Y_{1},\ldots,Y_{n}\) 是 \(S^{G}\) 中满足上面条件 (i),\ldots, (iv) 的元素。证明 \(S^{G}=K[Y_{1},\ldots,Y_{n}]\) 当且仅当
\end{enumerate}

\[
\operatorname{Card} (\mathrm{G}) = \frac {\prod_ {i} \deg \left(\mathrm{Y} _ {i}\right)}{\prod_ {i} \deg \left(\mathrm{X} _ {i}\right)}. *
\]

¶ 6) *设 n 是一个整数 ≥ 1，q 是一个素数的幂，K = F\_q，V = K\^{}n，G = GL(n, K)，G\_1 = SL(n, K)。把 G 与群 GL(V) 等同。此外，用 S 表示代数 S(V) = K{[}X\_1, \ldots, X\_n{]}，用 R（相应地，R\_1）表示 S 中由在 G（相应地，G\_1）下不变的元素组成的子代数。

\begin{enumerate}
\def\labelenumi{\alph{enumi})}
\tightlist
\item
  设 \(\mathbf{e} = (e_1, \ldots, e_n)\) 是一个整数序列 \(\geqslant 0\)。置
\end{enumerate}

\[
\mathrm{L} _ {\mathbf {e}} = \det (\mathrm{X} _ {i} ^ {q ^ {e _ {j}}}).
\]

这是 \(S\) 的一个元素。

证明

\[
g. \mathrm{L} _ {\mathbf {e}} = \det (g) \mathrm{L} _ {\mathbf {e}} \text {for all} g \in \mathrm{G}.
\]

（注意，若 \(g.X_i = \sum_j a_{ij} X_j\)，则 \(g.X_i^{q^e} = \sum_j a_{ij} X_j^{q^e}\)。）特别地，\(L_e\) 属于 \(R_1\)。

\begin{enumerate}
\def\labelenumi{\alph{enumi})}
\setcounter{enumi}{1}
\tightlist
\item
  设 \(j \in [1, n]\)。置
\end{enumerate}

\[
Z _ {j} = \prod_ {(a _ {i j})} (X _ {j} + \sum_ {i > j} a _ {i j} X _ {i}),
\]

其中乘积取遍 K 的元素的所有族 \((a_{ij})_{j < i\leqslant n}\)。设

\[
\mathrm{T} = \prod_ {j = 1} ^ {n} Z _ {j}.
\]

证明 T 整除所有 \(L_{e}\)。由此推出 \(T = L_{e_{n}}\)，其中 \(e_{n} = (0, 1, \ldots, n - 1)\)，且 T 在 \(G_{1}\) 下不变。

\begin{enumerate}
\def\labelenumi{\alph{enumi})}
\setcounter{enumi}{2}
\tightlist
\item
  用 \(\mathbf{e}_i\)（\(1 \leqslant i \leqslant n - 1\)）表示序列 \((0, 1, \ldots, i - 1, i + 1, \ldots, n)\)，并置 \(\mathrm{Y}_i = \mathrm{L}_{\mathbf{e}_i} / \mathrm{T}\)，cf.~b)。证明诸 \(\mathrm{Y}_i\) 属于 R 且 \(\deg(\mathrm{Y}_i) = q^n - q^i\)。
\end{enumerate}

\(d)\) 设 \(S' = K[X_1, \ldots, X_{n-1}]\)，并设 \(T', Y_1', \ldots, Y_{n-2}'\) 是 \(S'\) 中以与 \(T, Y_1, \ldots, Y_{n-1}\) 同样的方式（但把 \(n\) 换成 \(n-1\)）定义的元素。设 \(f: S \to S'\) 是由

\[
f \left(\mathrm{X} _ {i}\right) = \mathrm{X} _ {i} (1 \leqslant i \leqslant n - 1), f \left(\mathrm{X} _ {n}\right) = 0.
\]

定义的同态。证明

\[
f (\mathrm{T}) = 0, \quad f (\mathrm{Y} _ {1}) = \mathrm{T} ^ {\prime q (q - 1)}, \quad f (\mathrm{Y} _ {i}) = \mathrm{Y} _ {i - 1} ^ {\prime q} \quad \text{对} \quad 2 \leqslant i \leqslant n - 1.
\]

\begin{enumerate}
\def\labelenumi{\alph{enumi})}
\setcounter{enumi}{4}
\tightlist
\item
  证明族 \((T, Y_{1}, \ldots, Y_{n-1})\) 满足 Exerc. 5 的条件 (iv)。（设 \(x = (x_{1}, \ldots, x_{n})\) 是方程组 \((T, Y_{1}, \ldots, Y_{n-1})\) 在 K 的一个扩张 \(\overline{K}\) 中的一个零点。由于 x 是 T 的零点，诸 \(x_{i}\) 满足至少一个系数在 K 中的非平凡线性关系。必要时用 \(G_{1}\) 的一个元素变换 x，可以假设 \(x_{n} = 0\)。应用 d) 并对 n 用归纳法完成证明。）
\end{enumerate}

\(f)\) 证明 \(\mathbf{R}_1 = \mathbf{K}[\mathbf{T},\mathbf{Y}_1,\dots ,\mathbf{Y}_{n - 1}]\)，且 \(\mathrm{T},\mathrm{Y}_1,\dots ,\mathrm{Y}_{n - 1}\) 代数无关（“Dickson 定理”——应用 e) 与 Exerc. 5，并注意 \(\mathbf{G}_1\) 的阶等于 \(\deg (\mathrm{T}) = \sum_{i}\deg (\mathrm{Y}_i)\)）。

\begin{enumerate}
\def\labelenumi{\alph{enumi})}
\setcounter{enumi}{6}
\tightlist
\item
  证明 \(R = K[T^{q-1}, Y_{1}, \ldots, Y_{n-1}]\)。
\end{enumerate}

¶ 7) *设 R 是一个正则局部环（Commutative Algebra, Chap. VIII, §5, no. 1），其极大理想为 m，剩余域为 k。设 G 是 R 的一个有限自同构群，\(R' = R^{G}\) 是 R 中由在 G 下不变的元素组成的子环。这是一个局部环，其极大理想为 \(m' = m \cap R'\)。假设：

\begin{enumerate}
\def\labelenumi{(\roman{enumi})}
\item
  \(R'\) 是诺特的，且 R 是一个有限型 \(R'\)-模。
\item
  复合映射 \(R'\to R\to k\) 是满射。
\end{enumerate}

置 \(V = \mathfrak{m} / \mathfrak{m}^2\)；这是一个 \(k\)-向量空间。\(G\) 在 \(R\) 上的作用定义了一个同态 \(\varepsilon : G \to \mathbf{GL}(V)\)。

\begin{enumerate}
\def\labelenumi{\alph{enumi})}
\item
  设 \(\mathfrak{p}\) 是 \(R\) 的一个高度为 1 的素理想（Commutative Algebra, Chap. VII, §1, no. 6），\(s \in G\) 满足 \(s(\mathfrak{p}) = \mathfrak{p}\) 且 \(s\) 平凡地作用在 \(R / \mathfrak{p}\) 上。证明 \(\varepsilon(s)\) 是 \(V\) 的一个伪反射。（注意 \(\mathfrak{p}\) 在 \(\mathfrak{m} / \mathfrak{m}^2\) 中的像的维数是 0 或 1。）
\item
  证明：若 \(R'\) 正则，则子群 \(\varepsilon(G)\)（\(\mathbf{GL}(V)\) 的子群）由伪反射生成。（设 H 是 G 中由其在 \(\varepsilon\) 下的像为伪反射的元素生成的子群，\(R^{H}\) 是 R 中由在 H 下不变的元素组成的子环。利用 a) 证明：\(R'\) 的任何高度为 1 的素理想都不在 \(R^{H}\) 中分歧；利用 \(R^{H}\) 是整闭的这一事实，借助纯性定理 \(^{11}\) 推出 \(R' = R^{H}\)，从而 H = G。）
\item
  现在设 \(\mathbf{G}\) 的阶与 \(k\) 的特征互素。证明 \(\varepsilon\) 是单射。
\end{enumerate}

设 \((\mathfrak{m}_{n}^{\prime})\) 是 \(R^{\prime}\) 的、由 R 的 m-进滤链 \((\mathfrak{m}^{n})\) 诱导的滤链，并设

\[
i: \mathrm{gr} (\mathrm{R} ^ {\prime}) \to \mathrm{gr} (\mathrm{R})
\]

是从 \(R'\) 的相伴分次代数到 R 的相伴分次代数的典范同态（Commutative Algebra, Chap. III, §2）。证明 i 是单射，且它的像是子环 \(\mathrm{gr}(R)^{\mathrm{G}}\)，即 \(\mathrm{gr}(R)\) 中由在 G 下不变的元素组成的子环。

\(d)\) 沿用 \(c\) ) 的假设与记号，并进一步假设 \(\varepsilon (\mathbf{G})\) 由伪反射生成。若 \(l = \dim V\)，设 \(\mathrm{P}_1,\ldots ,\mathrm{P}_l\) 是 \(k\)-代数 \(\mathrm{gr}(R)^{\mathrm{G}}\) 的代数无关的齐次生成元（由 Th. 4 以及 \(\mathrm{gr}(R)\) 可与 V 的对称代数等同这一事实，这样的元素存在）；设 \(p_1,\ldots ,p_l\) 是它们的次数。由 \(c\) )，可以找到 \(x_{i}\in \mathfrak{m}_{p_{i}}^{\prime}\) 使得 \(\mathrm{gr}(x_i) = \mathrm{P}_i\)。证明诸 \(x_{i}\) 生成理想 \(\mathfrak{m}'\)，并由此推出 \(R'\) 是正则的。*

¶ 8) *设 V 是域 K 上一个有限维向量空间，S 是 V 的对称代数，G 是 GL(V) 的一个有限子群。假设代数 \(S^G\)（\(G\) 的对称不变量组成的代数）是一个分次多项式代数。

\begin{enumerate}
\def\labelenumi{\alph{enumi})}
\tightlist
\item
  设 u 是 V 的对偶 \(V^{*}\) 的一个元素，\(G_{u}\) 是 G 中由使 u 不变的元素组成的子群。证明 \(G_{u}\) 由伪反射生成。（线性型 u 可扩张为一个同态 \(f_{u}: S \to K\)。若用 \(S_{u}\) 表示 S 关于 \(f_{u}\) 的核的局部环，则环 \(S_{u}\) 是正则的。把 Exerc. 7 b) 应用于 \(G_{u}\)（把它视为 \(S_{u}\) 的自同构群），即得结论。）
\end{enumerate}

特别地，G 由伪反射生成。

\begin{enumerate}
\def\labelenumi{\alph{enumi})}
\setcounter{enumi}{1}
\tightlist
\item
  设 A 是 \(V^{*}\) 的一个子集，\(G_{A} = \bigcap_{u \in A} G_{u}\)。证明 \(G_{A}\) 由伪反射生成。（必要时扩张基域，可以假设 K 无限。证明在此情形下，存在由 A 生成的 \(V^{*}\) 的向量子空间中的一个元素 v，使得 \(G_{v} = G_{A}\)。对 \(v\) 应用 a) 即得结论。）
\end{enumerate}

\begin{enumerate}
\def\labelenumi{\arabic{enumi})}
\setcounter{enumi}{8}
\tightlist
\item
  设 V 是有限域 K（特征异于 2）上一个 4 维向量空间，Q 是 V 上一个指数为 2 的非退化二次型（Algebra, Chap. IX, §4, no. 2）。设 G = O(Q) 是这个型的正交群；这是一个由反射生成的有限群（loc. cit., §6, no. 4, Prop. 5）。
\end{enumerate}

\begin{enumerate}
\def\labelenumi{\alph{enumi})}
\item
  设 E 是 V 的一个极大全迷向子空间，\(G_{E}\) 是 G 中由使 \(g(x) = x\)（对所有 \(x \in E\)）的元素 g 组成的子群。证明 \(G_{E}\) 同构于 K 的加法群，且不含任何伪反射。
\item
  证明 G 的对称不变元代数不是一个分次多项式代数（利用前一习题）。
\end{enumerate}

